\chapter*{ВВЕДЕНИЕ}
\addcontentsline{toc}{chapter}{ВВЕДЕНИЕ}

Рекомендательные системы в настоящее
время играют значимую роль во многих сферах, связанных
с обработкой больших объемов информации, например
онлайн-магазины и кинотеатры, сервисы для прослушивания музыки.
Постоянный рост объема данных о поведении пользователей, их предпочтениях и интересах
требует использования инновационных решений для обработки, анализа и персонализации информации.
Это стало одной из причин увеличения роли рекомендательных систем~\cite{obzor_methods}.

Целью работы является проведение анализа методов построения рекомендательных систем.

Чтобы достигнуть поставленной цели, необходимо решить следующие задачи:
\begin{itemize}
\item провести анализ предметной области рекомендательных систем;
\item формализовать задачу составления рекомендаций;
\item провести обзор существующих методов или групп методов решения;
\item сформулировать критерии сравнения алгоритмов рекомендательных систем;
\item сравнить перечисленные методы по сформулированным критериям.
\end{itemize}

